On se place dans un repère $\big(O;~\vec{\imath},~\vec{\jmath}\big)$ orthogonal.

\begin{enumerate}
\item On considère la fonction $g$ définie pour tout réel $x$ par $g(x) = x^2 - 5x + 4$.

On note $P$ la courbe représentative de la fonction $g$.
	\begin{enumerate}
	\item Étudier le signe de la fonction $g$ sur $\mathbb{R}$.
	\item On considère un entier naturel $n$ quelconque.
	
	On note $A_n$ le point de la courbe $P$ d'abscisse $n$.
	
	On note $a_n$ le coefficient directeur de la droite $(A_nA_{n+1})$.

	Justifier que pour tout entier naturel $n$, on a $a_n = 2n - 4$.
	\item Quelle est la nature de la suite $(a_n)$ ?
	\end{enumerate}

\bigskip

\item On considère la fonction $f$ définie pour tout réel $x$ de l'intervalle $[0,5~;~8]$ par
\[f(x) = x - 5 + \dfrac{4}{x}.\]
On note $C$ la courbe représentative de la fonction $f$.
	\begin{enumerate}
	\item Vérifier que pour tout réel $x$ de l'intervalle $[0,5~;~8]$, on a $f(x) = \dfrac{g(x)}{x}$.
	\item À l'aide de la question \textbf{1.a}, déterminer la position de la courbe $C$ par rapport à l'axe des abscisses.
	\item On admet que la fonction $f$ est dérivable sur l'intervalle $[0,5~;~8]$.
	
	Montrer que pour tout réel $x$ de l'intervalle $[0,5~;~8]$, on a :
\[f'(x) = \dfrac{(x - 2)(x + 2)}{x^2}.\]
	\item En déduire le tableau de variations de la fonction $f$ sur l'intervalle $[0,5~;~8]$.
	\item Réaliser un schéma de l'allure de la courbe $C$ sur lequel apparaîtront les résultats des questions \textbf{2.b} et \textbf{2.d}.
	\end{enumerate}
\end{enumerate}
